% Options for packages loaded elsewhere
\PassOptionsToPackage{unicode}{hyperref}
\PassOptionsToPackage{hyphens}{url}
\documentclass[
  11pt,
]{article}
\usepackage{xcolor}
\usepackage[margin=1in]{geometry}
\usepackage{amsmath,amssymb}
\setcounter{secnumdepth}{-\maxdimen} % remove section numbering
\usepackage{iftex}
\ifPDFTeX
  \usepackage[T1]{fontenc}
  \usepackage[utf8]{inputenc}
  \usepackage{textcomp} % provide euro and other symbols
\else % if luatex or xetex
  \usepackage{unicode-math} % this also loads fontspec
  \defaultfontfeatures{Scale=MatchLowercase}
  \defaultfontfeatures[\rmfamily]{Ligatures=TeX,Scale=1}
\fi
\usepackage{lmodern}
\ifPDFTeX\else
  % xetex/luatex font selection
\fi
% Use upquote if available, for straight quotes in verbatim environments
\IfFileExists{upquote.sty}{\usepackage{upquote}}{}
\IfFileExists{microtype.sty}{% use microtype if available
  \usepackage[]{microtype}
  \UseMicrotypeSet[protrusion]{basicmath} % disable protrusion for tt fonts
}{}
\makeatletter
\@ifundefined{KOMAClassName}{% if non-KOMA class
  \IfFileExists{parskip.sty}{%
    \usepackage{parskip}
  }{% else
    \setlength{\parindent}{0pt}
    \setlength{\parskip}{6pt plus 2pt minus 1pt}}
}{% if KOMA class
  \KOMAoptions{parskip=half}}
\makeatother
\usepackage{graphicx}
\makeatletter
\newsavebox\pandoc@box
\newcommand*\pandocbounded[1]{% scales image to fit in text height/width
  \sbox\pandoc@box{#1}%
  \Gscale@div\@tempa{\textheight}{\dimexpr\ht\pandoc@box+\dp\pandoc@box\relax}%
  \Gscale@div\@tempb{\linewidth}{\wd\pandoc@box}%
  \ifdim\@tempb\p@<\@tempa\p@\let\@tempa\@tempb\fi% select the smaller of both
  \ifdim\@tempa\p@<\p@\scalebox{\@tempa}{\usebox\pandoc@box}%
  \else\usebox{\pandoc@box}%
  \fi%
}
% Set default figure placement to htbp
\def\fps@figure{htbp}
\makeatother
% definitions for citeproc citations
\NewDocumentCommand\citeproctext{}{}
\NewDocumentCommand\citeproc{mm}{%
  \begingroup\def\citeproctext{#2}\cite{#1}\endgroup}
\makeatletter
 % allow citations to break across lines
 \let\@cite@ofmt\@firstofone
 % avoid brackets around text for \cite:
 \def\@biblabel#1{}
 \def\@cite#1#2{{#1\if@tempswa , #2\fi}}
\makeatother
\newlength{\cslhangindent}
\setlength{\cslhangindent}{1.5em}
\newlength{\csllabelwidth}
\setlength{\csllabelwidth}{3em}
\newenvironment{CSLReferences}[2] % #1 hanging-indent, #2 entry-spacing
 {\begin{list}{}{%
  \setlength{\itemindent}{0pt}
  \setlength{\leftmargin}{0pt}
  \setlength{\parsep}{0pt}
  % turn on hanging indent if param 1 is 1
  \ifodd #1
   \setlength{\leftmargin}{\cslhangindent}
   \setlength{\itemindent}{-1\cslhangindent}
  \fi
  % set entry spacing
  \setlength{\itemsep}{#2\baselineskip}}}
 {\end{list}}
\usepackage{calc}
\newcommand{\CSLBlock}[1]{\hfill\break\parbox[t]{\linewidth}{\strut\ignorespaces#1\strut}}
\newcommand{\CSLLeftMargin}[1]{\parbox[t]{\csllabelwidth}{\strut#1\strut}}
\newcommand{\CSLRightInline}[1]{\parbox[t]{\linewidth - \csllabelwidth}{\strut#1\strut}}
\newcommand{\CSLIndent}[1]{\hspace{\cslhangindent}#1}
\setlength{\emergencystretch}{3em} % prevent overfull lines
\providecommand{\tightlist}{%
  \setlength{\itemsep}{0pt}\setlength{\parskip}{0pt}}
\usepackage{authblk}
\usepackage{helvet}
\usepackage{fontspec}
\setmainfont{Helvetica}
\usepackage[left]{lineno}
\usepackage{setspace}
\doublespacing
\usepackage{amsmath}
\usepackage{graphicx}
\usepackage{caption}
\usepackage{float}
\usepackage{siunitx}
\usepackage{placeins}
\usepackage{sectsty}
\sectionfont{\bfseries}
\subsectionfont{\bfseries}
\subsubsectionfont{\bfseries}
\usepackage[labelfont=bf]{caption}
\usepackage{xcolor}
\definecolor{captionblue}{HTML}{1F4E79}
\usepackage[labelfont={bf,color=captionblue}]{caption}
\usepackage{soul}
\sethlcolor{yellow}
\usepackage{hyperref}
\hypersetup{ colorlinks=false, pdfborder={0 0 0.5}, citebordercolor={0.8 0 0}, linkbordercolor={0 0 0.8}, urlbordercolor={0 0.5 0.8} }
\usepackage{xcolor}
\usepackage{bookmark}
\IfFileExists{xurl.sty}{\usepackage{xurl}}{} % add URL line breaks if available
\urlstyle{same}
\hypersetup{
  pdftitle={Mapping when genomic data can and cannot detect malaria transmission: biological constraints on inference},
  hidelinks,
  pdfcreator={LaTeX via pandoc}}

\title{\textbf{Mapping when genomic data can and cannot detect malaria
transmission: biological constraints on inference}}
\author{}
\date{\vspace{-2.5em}}

\begin{document}
\maketitle

\vspace{1mm}

\author{
Mouhamadou Fadel DIOP$^{1,2}$,
David Jeffries$^{1}$,
Nuredin Mohammed$^{1}$,
Martin Holland$^{4}$,
Alfred Amambua-Ngwa$^{1,2,3*}$
}

\noindent
$^{1}$ MRC Unit The Gambia at London School of Hygiene and Tropical Medicine, Banjul, The Gambia\\
$^{2}$ London School of Hygiene and Tropical Medicine, London, UK\\
$^{3}$ Department of Infection Biology, London School of Hygiene and Tropical Medicine, London, UK\\
$^{4}$ Wellcome Sanger Institute, Hinxton, UK
 
\vspace{5mm}
*Correspondence: \href{mailto:angwa@mrc.gm}{angwa@mrc.gm}; \href{mailto:alfred.ngwa@lshtm.ac.uk}{alfred.ngwa@lshtm.ac.uk}

\newpage

\section{Abstract}\label{abstract}

\textbf{Background.} Genomic surveillance is increasingly deployed to
track malaria transmission, yet the conditions under which genomic data
can reliably detect transmission links remain poorly defined. For
Plasmodium falciparum, the combined effects of high selfing rates,
variable recombination, and elevated background genomic relatedness may
fundamentally constrain inference---but this constraint has not been
systematically quantified across biologically relevant parameter spaces.

\textbf{Methods.} We developed a forward-time simulation framework using
SLiM, tskit, and msprime to generate full-genome sequence data under
three scenario classes: a baseline biological parameterisation, a
recombination sweep spanning four orders of magnitude (r = \(10^{-9}\)
to \(10^{-6}\) per base per generation), and a multi-population
structure with migration rates of 0.001--0.05. Transmission
detectability was evaluated using three inference approaches---IBD
(hmmIBD), identity-by-state (IBS), and maximum-likelihood phylogenetic
patristic distances (IQ-TREE)---across varying generation thresholds (G
= 1--25), sample sizes (n = 100--400), and population sizes. Performance
was quantified using area under the precision-recall curve (AUPRC), F1
score, and sensitivity at fixed specificity. Results were validated
against an observed P. falciparum dataset using same-household and
different-household pairs as a transmission proximity proxy.

\textbf{Results.} Transmission detectability was strongly and primarily
determined by recombination rate. AUPRC increased from 0.53--0.62 at r =
\(10^{-9}\) to 0.74--0.82 at r = \(10^{-6}\) regardless of migration
rate or sample size. IBD-based inference consistently outperformed IBS
and phylogenetic approaches, though the absolute advantage was small.
The generation threshold G was a critical secondary modulator: at G = 1,
no method exceeded AUPRC = 0.50 under any recombination condition.
Migration rate exerted minimal independent influence on detectability
across the range explored. In observed P. falciparum data, all three
methods performed near the random classifier baseline (AUPRC \textless{}
0.05), consistent with near-complete IBD distribution overlap between
household-linked and unlinked pairs.

\textbf{Conclusions.} Transmission detection from genomic data in P.
falciparum is subject to fundamental biological constraints that precede
methodological choices. Recombination rate and the temporal depth of the
transmission definition are the primary determinants of whether genomic
inference can succeed. Current IBD and IBS tools applied to endemic
African populations operate in a parameter regime where accurate
transmission detection is unlikely without integration of
epidemiological and spatial data. These results provide a quantitative
framework for calibrating expectations for genomic surveillance
programmes and identifying conditions under which genomic data alone is
insufficient to resolve transmission connectivity.

\textbf{Keywords}: Plasmodium falciparum; malaria genomic epidemiology;
IBD; transmission inference; simulation; recombination; detectability;
hmmIBD; IQ-TREE; AUPRC

\date{}
\maketitle
\linenumbers

\newpage

\section{\texorpdfstring{\textbf{Introduction}}{Introduction}}\label{introduction}

Understanding how malaria spreads between individuals and across
communities remains central to malaria control and elimination
(\citeproc{ref-WHO2023}{World Health Organization 2023}). Transmission
dynamics determine where interventions should be targeted, how rapidly
parasites move through populations, and whether local transmission
persists despite control efforts (\citeproc{ref-smith2010}{Smith et al.
2010}; \citeproc{ref-Griffin2010}{Griffin et al. 2010}). Over the past
decade, genomic data have become an increasingly integral component of
malaria surveillance, applied to distinguish local transmission from
importation, map parasite population connectivity, and reconstruct
transmission networks at fine spatial and temporal scales
(\citeproc{ref-Neafsey2015a}{Neafsey, Juraska, et al. 2015};
\citeproc{ref-Wesolowski2018}{Wesolowski et al. 2018};
\citeproc{ref-MalariaGEN2023}{MalariaGEN et al. 2021}). Large-scale
sequencing initiatives have demonstrated that genome-wide variation
captures relatedness gradients reflective of recent shared ancestry,
enabling probabilistic inference of epidemiological connectivity across
diverse endemic settings (\citeproc{ref-Daniels2015}{Daniels et al.
2015}; \citeproc{ref-Taylor2017}{Alyssa R. Taylor et al. 2017};
\citeproc{ref-Zhu2022}{Zhu et al. 2022};
\citeproc{ref-DiezBenavente2020}{Diez Benavente et al. 2020}).

Despite these advances, a critical and underexamined question persists:
under what biological and surveillance conditions can the genomic signal
of recent transmission actually be detected? Most genomic transmission
studies proceed under an implicit assumption that sufficient data
collection and appropriate analytical tools will recove meaningful
transmission signal (\citeproc{ref-Didelot2017}{Didelot et al. 2017};
\citeproc{ref-Campbell2018}{Campbell et al. 2018}). Yet P.
\emph{falciparum} poses a uniquely challenging inferential environment.
Obligate sexual recombination during reproduction in the mosquito midgut
progressively erodes haplotype-level relatedness between co-transmitted
parasite lineages, fragmenting shared ancestry into increasingly short
and sparse genomic segments with each successive transmission cycle
(\citeproc{ref-Conway1999}{Conway et al. 1999};
\citeproc{ref-Jiang2011}{Jiang et al. 2011};
\citeproc{ref-Miles2016}{Miles et al. 2016}). Polyclonal infections, in
which a single human host carries genetically distinct parasite clones,
further obscure pairwise genomic signals by introducing background
relatedness that is difficult to separate from transmission-derived
similarity (\citeproc{ref-Nkhoma2020}{Nkhoma et al. 2020};
\citeproc{ref-Felger2012}{Felger et al. 2012}). In endemic settings,
high background relatedness arising from population structure and
ongoing gene flow compounds these difficulties
(\citeproc{ref-ferreira2007}{Ferreira et al. 2007};
\citeproc{ref-chenet2012}{Chenet et al. 2012};
\citeproc{ref-AmambuaNgwa2019}{Amambua-Ngwa et al. 2019}). Together,
these features suggest that the limits of transmission signal detection
in malaria are not merely methodological, but biological.

The critical gap in the existing literature is the absence of a
principled characterisation of when the genomic signal of recent
transmission transitions from detectable to undetectable, independent of
the inference approach applied. Prior work has established that genomic
relatedness statistics can capture epidemiological connectivity under
favourable conditions (\citeproc{ref-Daniels2015}{Daniels et al. 2015};
\citeproc{ref-taylor2017}{Aimee R. Taylor et al. 2017};
\citeproc{ref-wong2017}{W. Wong et al. 2017};
\citeproc{ref-Henden2018}{Henden et al. 2018}), and that their
performance degrades under high recombination or complex population
structure (\citeproc{ref-schaffner2018}{Schaffner et al. 2018};
\citeproc{ref-Speed2015}{Speed and Balding 2015}). However, these
studies evaluate performance relative to inferred proxies or incomplete
empirical benchmarks, without access to known transmission ground truth,
making it difficult to determine whether inference failure reflects
signal absence or analytical inadequacy
(\citeproc{ref-Didelot2021}{Didelot et al. 2021};
\citeproc{ref-Albers2020}{Albers and McVean 2020}). As a result, the
parameter space in which transmission detection is biologically feasible
and the boundaries where it collapses regardless of analytical choice
remains uncharacterised.

Here, we develop a simulation-based framework to map the detectability
landscape of recent P. \emph{falciparum} transmission connectivity
across the key biological and surveillance dimensions that govern signal
availability. Using forward-time simulations with complete genealogical
ground truth (\citeproc{ref-haller2019}{Haller and Messer 2019};
\citeproc{ref-Kelleher2018}{Kelleher et al. 2018}), we define
transmission-linked ancestry precisely and quantify how it is
transformed, preserved, degraded, or rendered indistinguishable from
background as recombination rate, population structure, sampling
intensity, and transmission window depth vary. Rather than comparing
inference methods as competing approaches, we deploy three established
genomic summary statistics, identity by descent estimated via hmmIBD
(\citeproc{ref-schaffner2018}{Schaffner et al. 2018}), identity by state
(IBS), and phylogenetic inference from maximum-likelihood trees
(\citeproc{ref-Wong2025}{T. Wong et al. 2025}), as complementary
instruments for probing the boundaries of detectable signal. Consensus
failure across all three instruments is treated as evidence of a
biological detectability limit, not a methodological one
(\citeproc{ref-baumdicker2022}{Baumdicker et al. 2022};
\citeproc{ref-Speed2015}{Speed and Balding 2015}).

By mapping when genomic data can and cannot recover evidence of recent
transmission connectivity, this study reorients malaria genomic
epidemiology around a question that precedes the choice of analytical
method: does the signal exist? The boundaries delineated here define
realistic expectations for genomic transmission surveillance, clarify
the conditions under which investment in sequencing and analysis is
likely to return informative results, and provide a biologically
grounded framework for interpreting relatedness patterns in
field-collected parasite genomes. Rather than proposing a new inference
method, we establish the biological parameter space within which any
method can succeed and quantify how rapidly that space contracts as
recombination, sampling constraints, and transmission complexity
increase.

\subsection{A staged simulation framework for quantifying genomic
identifiability limits under controlled biological and surveillance
variation}\label{a-staged-simulation-framework-for-quantifying-genomic-identifiability-limits-under-controlled-biological-and-surveillance-variation}

In malaria, transmission occurs through \emph{Anopheles} mosquitoes
whose movements and infection histories are not individually tracked.
This vector-mediated biology means that genomic ancestry observed in
human infections reflects population-level mixing rather than discrete
person-to-person transmission chains, rendering direct reconstruction of
``who infected whom'' from sequence and/or epidemiological data largely
intractable. Rather than attempting to recover individual transmission
events, the analyses that follow quantify the conditions under which
genomic data can accurately recover transmission linked ancestry across
varying biological and surveillance constraints. To this end, we
constructed a simulation framework coupling three scenario blocks, a
baseline biological regime, a recombination rate sweep, and a structured
multi-population demography with a staged genomic simulation pipeline
and a multi-method inference layer evaluated against ground-truth IBD
using precision, recall, F1 and area under the precision-recall curve
(AUPRC) (\textbf{Figure 1}). This design ensures that observed limits on
recoverability are attributable to biological and demographic parameters
rather than to differences in analytical procedure. Simulation
parameters for all three scenarios are summarized in
\textbf{Supplementary Table S1}. The bottleneck effect on inference
performance is characterized separately in \textbf{Supplementary Figure
S1}, which shows that variation in effective population size at the
demographic bottleneck (Ne = 1,000 to 20,000) produces AUPRC differences
of 0.011--0.014 across methods at medium recombination, within the
confidence interval range and substantially smaller than the effects
attributable to recombination rate variation.

\vspace{5mm}

\pandocbounded{\includegraphics[keepaspectratio]{../results/main/figure1_framework.png}}

\emph{\textbf{Figure 1}. Overview of the simulation scenarios, genomic
data generation, and inference evaluation framework used for
transmission evaluation. The figure summarizes the end-to-end framework
linking simulated epidemiological scenarios to genomic inference and
performance evaluation. The figure outlines the conceptual and
computational structure used to evaluate limits of genomic transmission
identifiability. The left column defines the scenario space explored,
including a baseline biological setting, a recombination-rate sweep, and
structured multi-population models with migration. The central column
depicts the simulation pipeline, from forward-time simulations with
pedigree and tree-sequence recording, through tree-sequence processing
and mutation overlay, to generation of genome-wide variant data. The
right column summarizes the inference layer, distinguishing between
ground-truth ancestry (true identity by descent (IBD)), genomic
summaries derived from observable data (inferred IBD, identity by state
(IBS), and phylogenetic cophenetic distances from maximum-likelihood
trees), and downstream evaluation framework based on precision, recall,
and F1 metrics.}

\newpage

\section{\texorpdfstring{\textbf{Results}}{Results}}\label{results}

\subsection{\texorpdfstring{\textbf{1. Transmission-linked parasite
pairs retain only partial genealogical separation, establishing a
biological ceiling on detectable
signal.}}{1. Transmission-linked parasite pairs retain only partial genealogical separation, establishing a biological ceiling on detectable signal.}}\label{transmission-linked-parasite-pairs-retain-only-partial-genealogical-separation-establishing-a-biological-ceiling-on-detectable-signal.}

Before any inference layer is applied, the most fundamental question for
transmission detectability is whether the biological system produces
meaningfully distinct genealogical signals between directly linked and
unrelated parasite pairs. We addressed this by examining the
distribution of true coalescence times (TMRCA) between
transmission-linked pairs (direct descendants within a defined
generational window, G ≤ 25) and epidemiologically unrelated pairs
(non-direct) under baseline simulation parameters.

A prerequisite for any genomic method to distinguish transmission-linked
pairs from unlinked pairs is that the two classes occupy separable
regions of the genomic signal space. We examined this prerequisite using
the distribution of true pairwise coalescence times (TMRCA) under the
baseline and recombination sweep scenarios (\textbf{Figure 2}). In both
scenarios, direct transmission pairs (orange) display a distribution
shifted toward shorter TMRCA values relative to non-direct pairs (blue),
consistent with more recent common ancestry. However, the distributions
overlap considerably across the full range of TMRCA values observed,
with non-direct pairs contributing substantial density at low TMRCA
values in both scenarios.

Under the baseline scenario, direct pairs cluster below approximately 50
generations, but the tail of the non-direct distribution extends well
into this range, producing a region of irreducible overlap. This overlap
is at least as pronounced in the recombination sweep scenario, where
both distributions appear broader and their modes more similar. The
degree of overlap in TMRCA space does not markedly increase with
elevated recombination rate alone, suggesting that the biological
challenge of separation is present even before accounting for the
additional distortion introduced by the inference process. Together,
these distributions establish that the genomic signal separating direct
from non-direct transmission is inherently limited in this pathogen
system, and that any classifier operating on this signal will face a
non-trivial false positive burden regardless of method choice.

\vspace{5mm}

\pandocbounded{\includegraphics[keepaspectratio]{../results/main/figure2_tmrca_distribution.png}}

\emph{\textbf{Figure 2:} True TMRCA distributions for direct (orange)
and non-direct (blue) transmission pairs under the baseline scenario
(left panel) and recombination sweep scenario (right panel). The x-axis
shows true TMRCA in generations; the y-axis shows density. Each panel
pools all replicate simulations under the corresponding scenario.
Distributions are shown as kernel density estimates.}

\subsection{\texorpdfstring{\textbf{2. Recombination progressively
erodes the genomic signal of shared ancestry, widening the gap between
available and recoverable transmission
events.}}{2. Recombination progressively erodes the genomic signal of shared ancestry, widening the gap between available and recoverable transmission events.}}\label{recombination-progressively-erodes-the-genomic-signal-of-shared-ancestry-widening-the-gap-between-available-and-recoverable-transmission-events.}

We quantified how recombination rate modulates both the true genomic
relatedness signal and the capacity of hmmIBD to recover it. Across a
four-order-of-magnitude range of recombination rates (r = \(10^{-9}\) to
\(10^{-6}\) per base per generation), true pairwise IBD proportions
declined monotonically with increasing recombination (\textbf{Figure
3A}). At r = \(10^{-9}\), the distribution of true IBD is bimodal, with
a substantial fraction of pairs exhibiting IBD proportions approaching
1.0, consistent with near-complete identity in very low-recombination
genomes (\textbf{Figure S1}). As recombination rate increases to r =
\(10^{-6}\), this high-IBD mode collapses and the distribution
concentrates at lower values, reflecting the progressive fragmentation
of identity blocks by recombination.

The inferred IBD distribution from hmmIBD (\textbf{Figure 3B}) tracks
the true IBD distribution qualitatively across recombination rates, but
with systematic compression at intermediate values. The mean inferred
IBD proportion parallels the decay of true IBD, though consistently
underestimates it at low recombination rates and shows convergence at
higher rates (\textbf{Figure 3C}). The correlation between true and
inferred IBD decreases substantially with increasing recombination
(Figure 3D), declining from approximately 0.9 at r = \(10^{-9}\) to
roughly 0.55--0.60 at r = \(10^{-6}\). This progressive deterioration in
correlation indicates that the informational fidelity of hmmIBD is
itself a function of the recombination landscape, becoming increasingly
unreliable as recombination rate rises. The implication is that
detectability of transmission is partly constrained upstream of any
classification step, by the quality of the IBD signal that can be
extracted from the genome.

\vspace{5mm}

\pandocbounded{\includegraphics[keepaspectratio]{../results/main/figure3_ibd_recombination.pdf}}

\emph{\textbf{Figure 3}: IBD signal decay and inference fidelity across
recombination rates. (A) Distribution of true pairwise IBD proportions
for recombination rates spanning r =} \(10^{-9}\) to r = \(10^{-6}\),
shown as violin plots. (B) Distribution of IBD proportions inferred by
hmmIBD across the same recombination rates. (C) Mean true (solid) and
mean inferred (dashed) IBD proportions as a function of recombination
rate on a log scale, with 95\% confidence intervals. (D) Pearson
correlation between true and hmmIBD-inferred pairwise IBD proportions as
a function of recombination rate (log scale).

\subsection{\texorpdfstring{\textbf{3. Transmission detection undergoes
a threshold transition governed by recombination rate, independent of
inference
approach.}}{3. Transmission detection undergoes a threshold transition governed by recombination rate, independent of inference approach.}}\label{transmission-detection-undergoes-a-threshold-transition-governed-by-recombination-rate-independent-of-inference-approach.}

Precision-recall profiles and F1 scores varied substantially across
inference methods (IBD, IBS, phylogenetic patristic distances),
scenarios (baseline, recombination sweep, migration), and recombination
rates (\textbf{Figure 4}). In the precision-recall space (\textbf{Figure
4A}), all three methods show a wide dispersion of operating points, with
no method achieving uniformly high precision and recall simultaneously
across scenarios. IBD tends to cluster at higher precision values, while
IBS and phylogenetic-based inference show broader scatter toward lower
recall. The migration scenario introduces the most heterogeneity, with
operating points spread across a larger region of the precision-recall
plane compared to the baseline and recombination sweep.

\pandocbounded{\includegraphics[keepaspectratio]{../results/main/figure4_f1_recombinations.pdf}}

Figure 4. Classification performance of IBD, IBS, and phylogenetic
(Phylo) inference methods across simulation scenarios. (A)
Precision-recall scatter plot for all scenario-method combinations. Each
point represents one replicate simulation. Colour indicates scenario
(baseline, recombination sweep, migration); symbol shape indicates
method. (B) F1 score as a function of recombination rate for each
method, stratified by scenario. Lines show mean F1; shaded regions show
95\% confidence intervals across replicates. The x-axis is on a log
scale.

F1 scores across recombination rates (\textbf{Figure 4B}) reveal a
non-monotonic relationship in several scenario-method combinations.
Under the recombination sweep, F1 scores for all methods tend to be
lower at the extremes of the recombination range and modestly higher at
intermediate values, though the differences are small in absolute terms.
The baseline and migration scenarios display distinct F1 trajectories,
with the migration scenario generally yielding lower scores. Across all
recombination rates, no method consistently exceeds an F1 of
approximately 0.6, and many conditions yield F1 scores below 0.4. The
method comparison across sensitivity and false discovery rate
(\textbf{Figure Sx}) confirms that IBD achieves the most favourable
sensitivity-FDR trade-off across scenarios, whereas IBS and phylogenetic
inference accumulate higher false discovery rates at comparable
sensitivity levels, particularly under the migration scenario.

The cophenetic (patristic) distance distributions for direct and
non-direct pairs (\textbf{Figure 5}) provide additional context for the
phylogenetic method's performance. In both the baseline and
recombination sweep scenarios, the distributions of cophenetic distances
for direct and non-direct pairs overlap broadly. The degree of
separation between the two classes is modest in the baseline scenario
and does not improve in the recombination sweep, where elevated
recombination reshuffles ancestry and blurs phylogenetic signal. The
limited discriminative capacity visible in the distance distributions is
broadly consistent with the phylogenetic method's relatively low F1 and
precision-recall performance observed in \textbf{Figure 4}.

\vspace{5mm}

\pandocbounded{\includegraphics[keepaspectratio]{../results/main/figure5_phylogenetic_distances.pdf}}

\emph{\textbf{Figure 5:} Cophenetic (patristic) distance distributions
for direct (orange) and non-direct (blue) transmission pairs under the
baseline (left panel) and recombination sweep (right panel) scenarios.
Distances were derived from maximum-likelihood phylogenies inferred with
IQ-TREE. The x-axis shows cophenetic distance; the y-axis shows density.
Distributions are shown as kernel density estimates pooled across
replicates.}

\subsection{\texorpdfstring{\textbf{4. Migration rate has a limited
independent effect on detectability relative to
recombination.}}{4. Migration rate has a limited independent effect on detectability relative to recombination.}}\label{migration-rate-has-a-limited-independent-effect-on-detectability-relative-to-recombination.}

To disentangle the contributions of population structure and
recombination to transmission detectability, we evaluated AUPRC across a
factorial combination of recombination rates (\(10^{-9}\) to
\(10^{-6}\)) and migration rates (0.001, 0.01, 0.05) for G ≤ 25
generations (\textbf{Figure 6}). Across all panels (\textbf{Figure
6A--D}), AUPRC increased with recombination rate for all three methods,
regardless of migration level. IBD consistently achieved the highest
AUPRC among the three methods at every recombination-migration
combination tested. The gradient across migration rates within any fixed
recombination rate was modest: at r = \(10^{-6}\) (\textbf{Figure 6D}),
AUPRC for IBD varied by less than 0.05 units across the three migration
levels. The AUPRC gain curve (\textbf{Figure 6E}) confirms that the
primary driver of detectability differences across the panel is
recombination rate, not migration.

At fixed recombination rates, the sensitivity at 90\% specificity
(\textbf{Figure S3}) showed a similar pattern: IBD achieved the highest
sensitivity across all migration-recombination combinations, but
migration rate did not produce consistent directional changes in
sensitivity within any recombination stratum. The AUPRC profiles
stratified by migration rate (\textbf{Figure S6}) further demonstrate
that, at the AUPRC = 0.6 reference threshold, only the highest
recombination rate (r = \(10^{-6}\)) reliably produces AUPRC values
above this level for IBD, and only under some migration conditions. IBS
and Phylo remained below this threshold across nearly all conditions.
Sensitivity at 90\% specificity as a function of migration rate
(\textbf{Figure Sxxx}) shows that increasing migration from 0.001 to
0.05 does not substantially alter sensitivity in a consistent direction
for any method or recombination rate, suggesting that the range of
migration rates explored here represents a regime where connectivity
does not strongly modulate pair-level discriminability.

\vspace{5mm}

\pandocbounded{\includegraphics[keepaspectratio]{../results/main/Figure6_migration_redesigned.pdf}}

\emph{\textbf{Figure 6:} Mean AUPRC (± 95\% CI) for IBD (orange), IBS
(blue), and Phylo (green) methods across migration rates (0.001, 0.01,
0.05) and recombination rates (r =} \(10^{-9}\) to \(10^{-6}\)), for
pairs within G ≤ 25 generations. Panels A--D show AUPRC for each
recombination rate separately, with migration rate on the x-axis. Panel
E shows the AUPRC gradient across all recombination rates
simultaneously, with the three migration levels represented by line
colour. A horizontal reference line at AUPRC = 0.6 is included.

\subsection{\texorpdfstring{\textbf{5. Surveillance design shapes
detectability within the biological boundary but cannot transcend
it.}}{5. Surveillance design shapes detectability within the biological boundary but cannot transcend it.}}\label{surveillance-design-shapes-detectability-within-the-biological-boundary-but-cannot-transcend-it.}

To characterise the joint effect of recombination rate, sample size, and
migration on the upper bound of detectability, we summarised the best
AUPRC achievable by any single method (maximum over IBD, IBS, and Phylo)
across a grid of simulation conditions (\textbf{Figure 7}). The heatmaps
reveal a consistent gradient along the recombination axis: best AUPRC
increased from 0.53--0.62 at r = \(10^{-9}\) to 0.74--0.82 at r =
\(10^{-6}\), regardless of migration rate or sample size. This gradient
was present across all three migration panels (migration = 0.001, 0.01,
0.05) and was the dominant source of variation in the heatmaps.

The effect of sample size on best AUPRC was non-monotonic and generally
small relative to the recombination effect. At several
recombination-migration combinations, the smallest sample size (n = 100)
achieved the highest best AUPRC within its column for example, n = 100
yielded best AUPRC of 0.82 at r = \(10^{-6}\), migration = 0.001,
compared to 0.75 at n = 400 under the same conditions. This pattern was
not universal and was less consistent at lower recombination rates.
Across migration levels, the best AUPRC values were broadly similar: the
difference between migration = 0.001 and migration = 0.05 at any fixed
recombination-sample size cell was generally less than 0.05 AUPRC units.
These results reinforce the conclusion that recombination rate is the
primary axis of variation in transmission detectability within the
parameter space explored here, while sample size and migration exert
secondary and context-dependent effects.

Across the entire grid, best AUPRC did not reach 0.9 in any cell, and
only a small subset of conditions, specifically those combining the
highest recombination rate with smaller sample sizes, exceeded 0.8. This
ceiling on best achievable performance, even under the most favourable
simulated conditions, indicates a fundamental limit to transmission
detectability that persists despite method optimisation.

\pandocbounded{\includegraphics[keepaspectratio]{../results/main/Figure7_detectability.pdf}}

\emph{\textbf{Figure 7}: Best AUPRC (maximum over IBD, IBS, and Phylo
methods) across sample sizes (y-axis: 100, 200, 400) and recombination
rates (x-axis:} \(10^{-9}\), \(10^{-8}\), \(10^{-7}\), \(10^{-6}\)), for
three migration rates (panels: 0.001, 0.01, 0.05). Colour scale ranges
from yellow (low AUPRC ≈ 0.50) to dark red (high AUPRC ≈ 1.00). Numeric
AUPRC values are displayed in each cell. Results are for G ≤ 25
generations.

\subsection{\texorpdfstring{\textbf{6. The generation threshold for
transmission definition is a critical and underappreciated modulator of
detectability.}}{6. The generation threshold for transmission definition is a critical and underappreciated modulator of detectability.}}\label{the-generation-threshold-for-transmission-definition-is-a-critical-and-underappreciated-modulator-of-detectability.}

The definition of what constitutes a ``transmission link'' depends on
the temporal scope specified by the generation threshold G. We
systematically varied G from 1 to 25 generations and evaluated AUPRC for
each method across all recombination rates (\textbf{Figure 8}). At G =
1, the most stringent definition of transmission, AUPRC values were low
and nearly uniform across recombination rates for all three methods,
ranging from approximately 0.39 to 0.46 (\textbf{Figure 8A--C}). At this
threshold, no method achieved AUPRC above 0.50 under any recombination
condition, indicating that genomic inference of the most recent
transmission events is near-random in terms of ranking performance.

As G increased, AUPRC improved for all methods, but the rate of gain was
strongly modulated by recombination rate. At r = \(10^{-6}\), AUPRC for
IBD increased from 0.41 at G = 1 to 0.71 at G = 25, a gain of 0.30 AUPRC
units (\textbf{Figure 8A}). The corresponding gain for IBS was from 0.39
to 0.62 (\textbf{Figure 8B}), and for Phylo from 0.42 to 0.60
(\textbf{Figure 8C}). By contrast, at r = \(10^{-9}\), AUPRC increases
with G were smaller and less consistent across methods, IBD increased
from 0.42 to 0.49, while IBS and Phylo showed negligible gains. The
G-threshold gain curve (\textbf{Figure 8D}) illustrates this
recombination-G interaction: the curve is steepest at the highest
recombination rates, while at the lowest recombination rate, the AUPRC
gain with G is nearly flat.

The G-threshold heatmaps for all three methods (\textbf{Figure S7})
confirm these patterns at the replicate level, showing that the IBD
method consistently achieved the highest AUPRC at each G-recombination
combination, with the largest absolute advantage over IBS and Phylo at G
= 25 and high recombination. These results indicate that the
detectability of transmission events is fundamentally conditional on how
broadly ``transmission'' is defined temporally. Broader definitions
(larger G) yield more detectable signals, particularly in genomes with
higher recombination, whereas narrow definitions, most relevant for
outbreak or household-level epidemiology, remain poorly recovered by
current methods across the full recombination range explored.

\pandocbounded{\includegraphics[keepaspectratio]{../results/main/Figure8_G_threshold_interaction.pdf}}

\emph{\textbf{Fig 8:} AUPRC as a function of generation threshold G
(y-axis: G = 1, 3, 5, 10, 15, 25) and recombination rate (x-axis:}
\(10^{-9}\), \(10^{-8}\), \(10^{-7}\), \(10^{-6}\)), for IBD (panel A),
IBS (panel B), and Phylo (panel C). Numeric AUPRC values are shown in
each cell; colour scale from light yellow (low) to dark orange/red
(high). Panel D shows the G-threshold gain curve: mean AUPRC as a
function of G threshold, with one line per recombination rate and line
type per method (IBD, IBS, Phylo).

\subsection{\texorpdfstring{\textbf{7. The detectability landscape
defines a constrained but navigable parameter space for malaria genomic
surveillance.}}{7. The detectability landscape defines a constrained but navigable parameter space for malaria genomic surveillance.}}\label{the-detectability-landscape-defines-a-constrained-but-navigable-parameter-space-for-malaria-genomic-surveillance.}

To contextualise simulated performance against real sequencing data, we
evaluated IBD, IBS, and phylogenetic inference using an observed P.
\emph{falciparum} dataset in which same household (same-HH) and
different household (different-HH) pairs served as a proxy for
transmission proximity. Observed AUPRC values for all three methods
(\textbf{Figure 9B}) were low, approximately 0.03--0.05, placing them
far below the simulated performance distributions (mean simulated AUPRC:
\textasciitilde0.55--0.65 for IBD, \textasciitilde0.40--0.50 for IBS,
\textasciitilde0.25--0.35 for Phylo). The observed data points fall near
the dashed vertical line marking a random classifier (AUPRC ≈ prevalence
of positive pairs), indicating that none of the three methods could
meaningfully distinguish same-HH from different-HH pairs in this
dataset.

The IBD proportion distributions for same-HH (n = 158) and different-HH
(n = 12,562) pairs (\textbf{Figure 9C}) provide a direct explanation for
this failure: both distributions are nearly superimposable, with peaks
near IBD proportion ≈ 0.22--0.25. Same-HH pairs show a marginally
heavier tail at higher IBD proportions, but the overlap is extensive and
the distributions cannot be separated at any reasonable threshold. The
dashed vertical reference line in panel C indicates that the bulk of
both populations clusters within a narrow range of IBD values,
reflecting a population where background relatedness is high and
within-household transmission pairs are not detectably more related than
random pairs at the scale of pairwise IBD.

Overlaying the observed data AUPRC on the simulated heatmap
(\textbf{Figure 9A}) shows that the observed performance falls within
the region of the parameter space corresponding to low recombination
rates and the associated low-detectability zone. The dashed contour
lines in the migration = 0.001 and migration = 0.05 panels indicate that
the observed data are broadly consistent with simulations in the r =
\(10^{-7}\) to r = \(10^{-6}\) range at moderate to large sample sizes
and show AUPRC greater than the chosen threshold of 0.8, though the
mismatch between observed and simulated AUPRC suggests that additional
factors such as high background relatedness, sample aggregation
structure, or the imprecision of household membership as a transmission
proxy, further reduce performance beyond what the simulation captures.

Sensitivity analysis of the bottleneck size (\textbf{Figure S2}) showed
that AUPRC varied modestly with the effective population size at the
bottleneck (Ne bottleneck = 1,000--20,000) for both low (r =
\(10^{-9}\)) and medium (r = \(10^{-8}\)) recombination rates. The
smallest bottleneck produced marginally higher AUPRC, particularly for
IBD at low recombination, consistent with a more recent coalescent and
stronger relatedness signal under severe demographic contraction.
However, the absolute differences across bottleneck sizes were small,
indicating that the parameterisation of the bottleneck does not
qualitatively alter the conclusions about detectability limits.

The fraction of simulation replicates in which transmission events were
classifiable as identifiable---defined as exceeding a high-performance
AUPRC threshold---was near zero for the large majority of conditions
(Figure S4). Only IBD at r = \(10^{-9}\) achieved non-trivial
identifiable fractions (0.70--0.95 depending on migration level), and
this fraction dropped sharply to below 0.15 at r = \(10^{-8}\). IBS and
Phylo methods yielded identifiable fractions of 0.00 across all
recombination rates and migration levels, underscoring the near-complete
failure of these methods to achieve high-confidence transmission
classification in this pathogen system. Finally, the AUPRC gap between
IBD and the next-best method (\textbf{Figure S5}) was largest at the
lowest recombination rate (gap = 0.158--0.176 across migration levels)
and smallest at r = \(10^{-7}\) (gap ≈ 0.049--0.050), indicating that
the three methods are most differentiated in the regime where the
genomic signal is strongest. Across all conditions, IBD retained a
positive advantage over IBS and Phylo, though this advantage provides
limited practical benefit when absolute AUPRC values remain below 0.6 in
most conditions.

\pandocbounded{\includegraphics[keepaspectratio]{../results/main/Figure9_validation.pdf}}

\textbf{Figure 9}: (A) Heatmap of best simulated AUPRC (maximum over
IBD, IBS, Phylo) as a function of recombination rate (x-axis, log scale)
and sample size (y-axis: 100--400), for three migration rates (panels:
0.001, 0.01, 0.05). Dashed contour lines indicate the region of the
parameter space most consistent with observed P. \emph{falciparum} data.
Colour scale: green (AUPRC ≈ 0.5) to orange-red (AUPRC ≈ 1.0). (B) AUPRC
for IBD, IBS, and Phylo in observed (solid dots) versus simulated
(shaded interval) data. Horizontal dashed line marks the random
classifier baseline. (C) IBD proportion density distributions for
same-household (orange; n = 158) and different-household (grey; n =
12,562) pairs in the observed dataset. The vertical dashed lines
indicate the median IBD proportion for each group.*

\newpage

\section{Discussion}\label{discussion}

This study set out to map when genomic data can and cannot detect
malaria transmission, and the answer is unambiguous: in most
biologically relevant conditions simulated here, reliable discrimination
of transmission-linked parasite pairs from unlinked ones remains beyond
the reach of current methods. AUPRC rarely exceeded 0.65 under any
combination of method, recombination rate, migration, or sample size,
and performance for Plasmodium \emph{falciparum} relevant parameters
particularly low recombination and high background IBD was substantially
worse, approaching the random classifier baseline in the observed
dataset. These findings are broadly consistent with the theoretical
expectation that parasites with high rates of selfing and low effective
recombination will maintain extensive identity by descent across
unrelated transmission chains, blurring the genomic boundary between
linked and unlinked pairs (\citeproc{ref-Neafsey2015b}{Neafsey,
Waterhouse, et al. 2015}; \citeproc{ref-Manske2012}{Manske et al.
2012}). However, the degree to which this boundary is irrecoverable even
with ideal IBD-based inference was not previously quantified in a
simulation framework that jointly varied recombination, population
structure, and temporal transmission depth.

The dominance of IBD-based inference over IBS and phylogenetic
approaches across essentially all conditions is consistent with prior
work demonstrating that haplotype-level sharing provides more precise
relatedness estimates than allele-by-allele identity for highly selfing
organisms (\citeproc{ref-Henden2018}{Henden et al. 2018};
\citeproc{ref-Apinjoh2019}{Apinjoh et al. 2019}). Schaffner et
al.~established hmmIBD as the reference tool for pairwise IBD estimation
in P. \emph{falciparum} and noted that inference quality degrades with
increasing recombination (\citeproc{ref-schaffner2018}{Schaffner et al.
2018}), a pattern directly replicated here in the declining correlation
between true and inferred IBD across \textbf{Figure 3}. Chang et
al.~similarly showed that genetic relatedness measures could distinguish
recent from older transmission in low-transmission settings, but their
analyses were conducted in populations with relatively sparse background
relatedness (\citeproc{ref-Auburn2018}{Auburn et al. 2018}). Our results
suggest that this finding is not readily generalisable to
high-transmission, high-IBD settings: when background relatedness is
elevated, as in the observed dataset (\textbf{Figure 9C}), even true IBD
provides insufficient discriminatory power, irrespective of inference
accuracy. This represents a biological, not methodological, ceiling on
detectability.

The finding that recombination rate was the dominant driver of
transmission detectability exceeding the effects of migration rate,
sample size, and G threshold in most conditions extends and refines the
theoretical framework established in population genetics of malaria.
Nkhoma et al.~demonstrated that recombination breaks down extended
haplotype structure in P. \emph{falciparum}, reducing IBD block length
and eroding genealogical signal over time
(\citeproc{ref-Nkhoma2012}{Nkhoma et al. 2012}). Our results place this
erosion in a quantitative performance framework: the transition from r =
\(10^{-9}\) to r = \(10^{-6}\) per base per generation translated to a
gain of 0.22--0.30 AUPRC units across methods (\textbf{Figure 8}), far
larger than any variation attributable to population structure or sample
size. This has practical consequences. Genomic epidemiology programmes
that benchmark tools using high-recombination reference datasets that
implicitly assume recombination rates comparable to other pathogens may
substantially overestimate detectability in the field. The WHO supported
adoption of genomic surveillance for malaria
(\citeproc{ref-Didelot2014}{Didelot, Gardy, and Colijn 2014}) needs to
be grounded in explicit estimates of local recombination parameters
before performance expectations can be set.

The role of the generation threshold G as a modulator of detectability
has not previously been explored systematically in the malaria genomics
literature, yet our results show it to be a first order determinant of
performance. At G = 1, operationally the most relevant definition for
outbreak or household transmission investigations, AUPRC was below 0.50
for all methods and all recombination rates. This is sobering. The
phylogenetic transmission inference literature in bacterial pathogens,
particularly tuberculosis and SARS-CoV-2, has broadly assumed that
phylogenetic proximity is a reasonable proxy for recent transmission
(\citeproc{ref-Didelot2014}{Didelot, Gardy, and Colijn 2014};
\citeproc{ref-DeMaio2015}{De Maio et al. 2015}). In those organisms,
mutation rates are orders of magnitude higher and recombination is
absent or infrequent, making temporal signal recoverable from sequence
data alone. In P. \emph{falciparum}, neither condition holds. The TMRCA
overlap between direct and non-direct pairs observed in \textbf{Figure
2} is a direct consequence of the organism's biology, and the
G-threshold analysis (\textbf{Figure 8}) confirms that the genomic
signal necessary to resolve very recent transmission is simply not
present at any recombination rate in the parameter space we explored, a
finding with no clear analogue in the bacterial or viral transmission
inference literature (\citeproc{ref-MohdAbdRazak2016}{Mohd Abd Razak et
al. 2016}).

Beyond benchmarking, these results carry implications for the
theoretical framing of genomic transmission inference as an
identifiability problem. In the formal sense, a statistical inference
problem is identifiable if distinct parameter values produce distinct
observable distributions. The TMRCA overlap in \textbf{Figure 2} and the
near-superimposable IBD distributions in \textbf{Figure 9C} together
suggest that, under realistic malaria transmission conditions, the
distributions of genomic distance between direct and non-direct pairs
are not reliably distinguishable. This constitutes a form of practical
non-identifiability that precedes any modelling assumption about
transmission chains. Didelot et al.~formalised similar concerns in the
context of bacterial whole-genome sequence phylogenetics, arguing that
without strong epidemiological priors, transmission reconstruction is
severely underdetermined (\citeproc{ref-Didelot2020}{Didelot, Siveroni,
and Volz 2020}). The malaria case is arguably more extreme because the
genomic signal itself, not merely its interpretation, is contaminated by
background relatedness and short IBD blocks. Moving forward, frameworks
that integrate genomic data with spatial, temporal, and household-level
epidemiological data are likely to be necessary to approach reliable
transmission inference in high transmission settings (@
\citeproc{ref-Wesolowski2018}{Wesolowski et al. 2018};
\citeproc{ref-Hay2002}{Hay et al. 2002}).

\section{Limitations and scope}\label{limitations-and-scope}

Several limitations of this study should be considered when interpreting
these findings. First, simulations were restricted to a single
chromosome (chromosome 1, 640,851 bp) of the P. \emph{falciparum} genome
rather than the full 14 chromosomes, with substantial heterogeneity in
recombination rate, GC content, and functional constraint across the
genome. In particular, telomeric and subtelomeric regions harboring
hypervariable antigen gene families (var, rifin, stevor) violate the
neutral evolution assumptions of both the simulation model and hmmIBD,
and are typically excluded from relatedness analyses in practice.
Chromosomes also differ in effective length and recombination landscape,
which will produce chromosome-specific AUPRC values around the
genome-wide average reported here. While chromosome 1 provides a
tractable proof-of-concept framework, IBD statistics computed from the
whole genome are expected to be more precise due to greater averaging
over independent chromosomal segments. Whole-genome simulations would
likely increase inferred IBD accuracy and may modestly raise the
detectability ceiling, though the biological signal constraints
identified here, particularly TMRCA overlap, would remain.

Second, the model does not incorporate within-host diversity (i.e.,
complexity of infection or multiplicity of infection, COI/MOI). Although
COI was not explicitly modelled, the recombination sweep scenario
partially captures its genomic consequences, as higher recombination
rates mimic the increased haplotype fragmentation and reduced IBD block
length that arise when multiple co-infecting clones contribute to the
genomic background of a sampled parasite. Mixed infections are common in
high-transmission settings and introduce substantial uncertainty into
pairwise relatedness estimates, as hmmIBD and IBS-based tools perform
less reliably on polyclonal samples {[}Henden et al.
(\citeproc{ref-Henden2018}{2018}); 15{]}. Our simulations therefore
likely overestimate achievable field performance, as they implicitly
assume monoclonal, correctly phased haplotypes. Third, the bottleneck
parameterisation, while varied in sensitivity analysis (\textbf{Figure
S2}), does not capture the full complexity of demographic history in
endemic populations, including recent range expansion, seasonal
bottlenecks, or the impact of malaria control interventions on Ne
(\citeproc{ref-AmambuaNgwa2019}{Amambua-Ngwa et al. 2019}). Fourth, the
study lacks a formal identifiability analysis in the mathematical sense;
the empirical AUPRC approach quantifies discriminative performance but
does not derive analytical bounds on what is recoverable from the data.
Finally, there is a potential circularity between the simulation model
and inference: hmmIBD was designed with assumptions about recombination
and selfing rates that partially overlap with those used to generate the
simulation data, which may produce modestly optimistic IBD inference
performance relative to a truly independent generative model.

\section{Future Research Directions}\label{future-research-directions}

Several directions follow directly from these findings. Whole-genome
simulations incorporating all 14 P. \emph{falciparum} chromosomes,
including chromosome-specific recombination rates and the hypervariable
subtelomeric regions, are a necessary next step to assess whether
detectability limits identified here are specific to the
single-chromosome setting or hold genome-wide. Incorporating polyclonal
infection models with realistic COI distributions would substantially
improve the biological fidelity of the simulation framework and is
likely to reveal additional performance losses not captured here.
Genome-wide simulation extending across all 14 chromosomes with
empirical recombination maps (Miles et al.~2016; MalariaGEN) represents
an important direction for future work, and would permit direct
calibration of the recombination-driven identifiability ceiling reported
here against the full genomic architecture of P. \emph{falciparum}.
Methods that jointly model genomic relatedness, spatial distance, and
household structure, as demonstrated in principle in some bacterial and
viral settings (\citeproc{ref-Wesolowski2018}{Wesolowski et al. 2018};
\citeproc{ref-Hay2002}{Hay et al. 2002}), should be adapted and
evaluated for P. \emph{falciparum}, with explicit attention to the
high-relatedness, low-recombination regime that characterises endemic
African populations. Empirical validation studies using datasets where
household or clinic-confirmed transmission chains are known with higher
certainty than household membership alone would also allow sharper
calibration of the detectability limits mapped here. Lastly, formalising
the non-identifiability conditions described in this paper, deriving
analytical bounds on the minimum recombination rate, sample size, or
temporal resolution required for transmission inference to achieve a
specified precision, would provide the theoretical foundation currently
missing from the malaria genomic epidemiology literature and could
inform the design of future surveillance systems before costly
sequencing investments are made.

\section{Conclusions}\label{conclusions}

This study set out to define the conditions under which genomic data can
and cannot detect malaria transmission, using a simulation-based
framework that jointly varied recombination rate, population structure,
migration, sample size, and the temporal scope of the transmission
definition. The central finding is that transmission detectability in P.
falciparum is primarily constrained by recombination rate, with AUPRC
rising substantially across the four order of magnitude range explored,
and that no method, IBD, IBS, or phylogenetic inference, achieved
reliable discriminative performance in the low-recombination,
high-background relatedness regime most characteristic of endemic
African settings. The generation threshold G emerged as a critical
secondary determinant, with performance below chance for all methods at
G = 1, the operationally most relevant temporal scope for outbreak and
household transmission investigations. Validation against observed P.
falciparum data confirmed that current tools are unlikely to resolve
transmission links in high-transmission settings: the near-identical IBD
distributions between household-linked and unlinked pairs indicate a
biological signal ceiling that inference tools cannot overcome.

These findings reframe the problem of transmission inference in malaria
as one of biological constraint rather than methodological inadequacy.
The TMRCA overlap between direct and non-direct transmission pairs
observed in simulation is not an artefact of tool design but a
consequence of the organism's high selfing rate, high background
relatedness, and short IBD blocks under endemic transmission intensity.
In this sense, the study contributes a formal, quantitative dimension to
a concern that has been raised qualitatively in the genomic epidemiology
literature but not previously tested across the parameter space that
governs field performance. The results argue for a reorientation of
expectations: genomic data in malaria may be most informative not as a
standalone transmission discriminator, but as one layer within an
integrated inference framework that incorporates spatial proximity,
temporal co-occurrence, and household or clinical contact data.

As genomic surveillance programmes expand across sub-Saharan Africa and
Southeast Asia, the question of what parasite genomes can reasonably
tell us about transmission is no longer theoretical. This study provides
a quantitative answer for a defined parameter space and demonstrates
that the answer is sobering in most endemic conditions. Moving forward,
the framework developed here, simulation-based detectability mapping
across biological parameter gradients, offers a template for
pre-deployment evaluation of genomic surveillance tools in any pathogen
system where background relatedness, recombination, and population
structure jointly shape the feasibility of transmission inference.
Translating these insights into surveillance design, sampling
strategies, and analytical pipelines that are appropriately calibrated
for the local biological context will be essential to ensuring that
genomic investments in malaria control deliver actionable
epidemiological intelligence.

\newpage

\section{Methods}\label{methods}

\subsection{Overview of the simulation
framework}\label{overview-of-the-simulation-framework}

We developed a comprehensive simulation framework to evaluate the
accuracy of genomic methods for reconstructing Plasmodium
\textbf{falciparum} transmission networks. Our approach integrates
forward-time population genetic simulation using SLiM v4.0 (Haller \&
Messer, 2023) with coalescent-based genealogical reconstruction using
msprime v1.2 (Kelleher et al., 2016) and tree sequence recording via
pyslim v1.0 (Haller et al., 2019). This hybrid framework allows us to
capture both the selective dynamics of recent parasite evolution and the
deeper coalescent history while maintaining computational tractability
for genome-scale simulations.

We simulated three distinct scenarios to evaluate transmission inference
under varying evolutionary conditions:

\begin{enumerate}
\def\labelenumi{\arabic{enumi}.}
\tightlist
\item
  \textbf{Baseline scenario}: Fixed P. \textbf{falciparum} specific
  parameters representing typical endemic transmission
\item
  \textbf{Recombination sweep}: Systematic variation of recombination
  rates to assess the impact of within-host diversity on transmission
  reconstruction.
\item
  \textbf{Full parameter space}: Factorial exploration of recombination
  rate, transmission bottleneck size, mutation accumulation (expected
  substitutions per transmission, EST), sampling intensity, outbreak
  size, and population structure with migration.
\end{enumerate}

\subsection{Biological Model and
Assumptions}\label{biological-model-and-assumptions}

\subsubsection{Parasite Life Cycle and Reproductive
Mode}\label{parasite-life-cycle-and-reproductive-mode}

P. \textbf{falciparum} exhibits a complex haploid-diploid life cycle
alternating between human and mosquito hosts. Our simulations explicitly
model this biology with the following key assumptions:

\textbf{Haploidy}: The parasite genome remains haploid throughout
asexual replication in both human (liver and blood stages) and mosquito
(sporozoite) hosts. Diploidy occurs only transiently during sexual
reproduction in the mosquito midgut (zygote/ookinete stage), followed
immediately by meiosis. We therefore model parasites as fundamentally
haploid organisms (one genome copy per individual) with occasional
outcrossing.

\textbf{High selfing rate}: Following empirical estimates from endemic
populations (Daniels et al., 2015), we set the selfing rate
\(\sigma = 0.95\) for baseline simulations. This reflects the
predominantly clonal transmission of P. \textbf{falciparum} due to: low
multiplicity of infection (MOI) in many endemic settings, temporal
asynchrony of gametocyte production and stochastic sampling of
gametocytes during mosquito blood meals.

The effective outcrossing rate is thus \(\tau = 1 - \sigma = 0.05\).

\textbf{Discrete generations}: Each complete transmission cycle (human →
mosquito → human) represents one discrete generation in our model. This
comprises: - Human liver stage: 7-10 days (1 generation) - Human blood
stage: 5 asexual replication cycles over \textasciitilde10 days (5
generations) - Mosquito sexual stage: Gametocyte → ookinete (1-2
generations) - Mosquito asexual stage: Oocyst sporogony (8-10
generations) - Mosquito sporozoite migration: (1-2 generations)

The total generation time per transmission event is approximately
\(g_{\text{trans}} = 22\) generations (range: 18-28 depending on
infection duration), spanning 12-15 days of calendar time.

\subsubsection{Transmission Bottleneck}\label{transmission-bottleneck}

A critical feature of P. \textbf{falciparum} transmission is the severe
genetic bottleneck occurring during mosquito-to-human transmission. When
an infected mosquito takes a blood meal, only a small number of
sporozoites successfully invade hepatocytes and establish blood-stage
infection. We model this bottleneck explicitly with parameter \(B\),
representing the number of parasite genomes transmitted per infection
event.

We explore three bottleneck scenarios: \textbackslash{} - \textbf{Tight
bottleneck}: \(B = 1\) (single parasite transmission, severe genetic
drift) - \textbf{Moderate bottleneck}: \(B = 5\) (few parasites,
intermediate drift) - \textbf{Loose bottleneck}: \(B = 20\) (multiple
parasites, reduced drift)

Empirical estimates suggest \(B \approx 1-10\) in natural transmission
(Schneider et al., 2018), making this a key parameter governing the
correlation between genetic and epidemiological relationships.

\subsubsection{Population Structure and
Demography}\label{population-structure-and-demography}

\textbf{Single-population (baseline) model}: We simulated a single
panmictic population with an ancient effective population size:
\(N_e = 10,000\) (representing pre-bottleneck ancestral diversity) -
Recent effective population size: \(N_0 = 1,000\) (representing
contemporary population) - Exponential population size change occurring
at generation \(g_{\text{change}} = 200\) before present:

\[N(t) = N_0 \left(\frac{N_e}{N_0}\right)^{t/g_{\text{change}}}\]

where \(t\) represents generations before present
(\(0 \leq t \leq g_{\text{change}}\)).

\textbf{Multi-population (migration) model}: For the full parameter
space exploration, we simulated \(n = 3\) populations (demes) with
symmetric migration at rate \(m\) per generation. The migration matrix
\(\mathbf{M}\) is:

\[M_{ij} = \begin{cases}
1 - (n-1)m & \text{if } i = j \\
m & \text{if } i \neq j
\end{cases}\]

We explored three migration scenarios:

\begin{itemize}
\tightlist
\item
  Low gene flow: \(m = 0.001\) (1 migrant per 1000 generations)
\item
  Moderate gene flow: \(m = 0.01\) (1 migrant per 100 generations)
\item
  High gene flow: \(m = 0.05\) (1 migrant per 20 generations)
\end{itemize}

This represents varying degrees of spatial structure from nearly
isolated populations to high connectivity, reflecting different endemic
transmission settings (isolated villages vs.~urban transmission
hotspots).

\subsection{Evolutionary Parameters}\label{evolutionary-parameters}

\subsubsection{Recombination}\label{recombination}

Recombination in P. \textbf{falciparum} occurs during meiosis following
sexual reproduction in the mosquito midgut. The recombination rate per
base pair per generation, \(r\), determines the extent to which parental
genomes are reshuffled during outcrossing events.

We systematically varied \(r\) across four orders of magnitude:

\begin{itemize}
\tightlist
\item
  Very low: \(r = 10^{-9}\) bp\(^{-1}\) gen\(^{-1}\) (effectively
  clonal)
\item
  Low: \(r = 10^{-8}\) bp\(^{-1}\) gen\(^{-1}\) (rare recombination)
\item
  Medium: \(r = 10^{-7}\) bp\(^{-1}\) gen\(^{-1}\) (moderate
  recombination)
\item
  High: \(r = 10^{-6}\) bp\(^{-1}\) gen\(^{-1}\) (frequent
  recombination)
\end{itemize}

Because direct simulation of polyclonal infections was beyond the scope
of the current framework, effective recombination rate was used as a
proxy for within-host complexity: elevated recombination rates
approximate the increased outcrossing and haplotype mixing that occur
when multiple genetically distinct clones co-infect a single host,
thereby capturing, in aggregate, the signal erosion associated with
higher multiplicity of infection.

For a chromosome of length \(L = 640,851\) bp (chromosome 1), the
genome-wide recombination rate is \(R = 4N_e r L\), ranging from
effectively no recombination (\(R \approx 0.03\) for \(r = 10^{-9}\)) to
free recombination (\(R \approx 26\) for \(r = 10^{-6}\)).

\subsubsection{Mutation}\label{mutation}

Two mutation processes were modeled as follow:

\textbf{Neutral mutations}: Added post-simulation using msprime with
rate \(\mu = 10^{-8}\) bp\(^{-1}\) gen\(^{-1}\), consistent with
empirical estimates for P. \textbf{falciparum} (Hamilton et al., 2017).
These mutations accumulate according to the infinite-sites model and
provide genomic variation for IBD/IBS inference.

\textbf{Selected mutations}: To model adaptive evolution (drug
resistance), we introduced a single beneficial allele at position
\(\ell_s = L/3\) with: - Selection coefficient: \(s = 0.3\) (30\%
fitness advantage) - Dominance coefficient: \(h = 0.5\) (codominant,
though relevant only during transient diploid stage) - Introduction
time: \(g_s = 80\) generations before present

The selected allele was introduced into \(n_{\text{origins}} = 1\)
random genome and tracked until fixation or loss. Simulations where the
beneficial allele was lost were automatically restarted (maximum 100
restarts) to condition on successful selective sweeps, allowing us to
evaluate transmission inference in the presence of recent positive
selection.

\subsubsection{Expected Substitutions per Transmission
(EST)}\label{expected-substitutions-per-transmission-est}

A key parameter linking mutation rate to transmission dynamics is the
expected number of mutations accumulated per transmission event. For a
genome of length \(L\) transmitted over \(g_{\text{trans}}\) generations
with mutation rate \(\mu\):

\[\text{EST} = \mu \times L \times g_{\text{trans}}\]

We explored three EST scenarios:

\begin{itemize}
\tightlist
\item
  None: EST \(= 0\) (mutations added only post-simulation via msprime)
\item
  Low: EST \(= 0.1\) (rare mutations during transmission)
\item
  Moderate: EST \(= 0.5\) (detectable mutation accumulation)
\end{itemize}

For EST \(> 0\), we calculated the corresponding per-generation mutation
rate as:

\[\mu_{\text{trans}} = \frac{\text{EST}}{L \times g_{\text{trans}}}\]

Using \(L = 640,851\) bp and \(g_{\text{trans}} = 22\) generations
(typical scenario):

\[\mu_{\text{trans}} = \frac{\text{EST}}{640,851 \times 22} = \frac{\text{EST}}{14,098,722} \approx 7.1 \times 10^{-8} \times \text{EST}\]

For EST \(= 0.5\): \(\mu_{\text{trans}} \approx 3.5 \times 10^{-8}\)
bp\(^{-1}\) gen\(^{-1}\).

\subsection{Mathematical Framework}\label{mathematical-framework}

\subsubsection{Wright-Fisher Model with Selection and
Migration}\label{wright-fisher-model-with-selection-and-migration}

Our simulations implement a discrete-generation Wright-Fisher model
modified for haploid, selfing, and selection. The probability that
genome \(i\) in generation \(t+1\) descends from genome \(j\) in
generation \(t\) is:

\[P(i_{t+1} \leftarrow j_t) = \frac{w_j}{\bar{w}} \times \left[\sigma + (1-\sigma) \times \frac{1}{N}\right]\]

where \(w_j\) is the fitness of genome \(j\),
\(\bar{w} = \frac{1}{N}\sum_{k=1}^N w_k\) is the mean population
fitness, \(\sigma\) is the selfing rate and \(N\) is the population
size.

For the multi-population model, migration occurs before reproduction.
The probability that genome \(i\) in population \(p\) at generation
\(t+1\) descends from population \(q\) is \(M_{pq}\), and given
migration to population \(p\), the probability of descending from a
specific genome follows the within-population model above.

\subsubsection{Fitness Model}\label{fitness-model}

For a genome carrying \(n\) copies of the beneficial allele at the
selected locus (0 or 1 in haploids):

\[w = \begin{cases}
1 & \text{if } n = 0 \text{ (wild-type)} \\
1 + s & \text{if } n = 1 \text{ (mutant)}
\end{cases}\]

During transient diploidy (sexual stage), heterozygous genotypes have
fitness \(w = 1 + hs\), but this affects only the single generation of
meiosis. The derived allele frequency in generation \(t\), \(p_t\),
evolves deterministically (ignoring drift) as:

\[p_{t+1} = p_t \times \frac{1+s}{1 + s p_t}\]

With drift in a finite population, frequency changes are stochastic,
sampled from a binomial distribution:

\[X_{t+1} \sim \text{Binomial}\left(N, p_t \times \frac{1+s}{1 + s p_t}\right)\]

where \(X_t\) is the count of derived alleles.

\subsubsection{Transmission Bottleneck
Model}\label{transmission-bottleneck-model}

Each transmission event samples \(B\) parasite genomes from the donor's
infection. For a donor with derived allele frequency
\(p_{\text{donor}}\), the probability that the recipient infection has
\(k\) copies of the derived allele is:

\[P(k \mid p_{\text{donor}}, B) = \binom{B}{k} p_{\text{donor}}^k (1-p_{\text{donor}})^{B-k}\]

For tight bottlenecks (\(B=1\)), transmission is deterministic:
\textbackslash{}

\(P(k=1) = p_{\text{donor}}\), \(P(k=0) = 1-p_{\text{donor}}\)

This bottleneck severely constrains the genetic diversity transmitted
and creates strong stochasticity in allele frequency trajectories.

\subsubsection{Coalescent Theory and Tree
Sequences}\label{coalescent-theory-and-tree-sequences}

The genealogical relationships among sampled parasites are recorded as
tree sequences (Kelleher et al., 2018), encoding the full ancestral
recombination graph (ARG). For neutral loci in a panmictic population of
size \(N_e\), the distribution of pairwise coalescence times \(T\)
follows:

\[P(T > t) = e^{-t/(2N_e)}\]

with mean coalescence time \(\mathbb{E}[T] = 2N_e\) generations.

Under the sequentially Markovian coalescent (SMC) approximation,
genealogies at different genomic positions are correlated, with the
correlation decaying as a function of recombination distance \(d\):

\[\rho(d) = e^{-\rho d}\]

where \(\rho = 4N_e r\) is the population-scaled recombination rate.

Tree sequences enable efficient computation of true \textbf{Identity by
descent (IBD)} metric, which is the genomic segments sharing a common
ancestor within \(T_{\text{IBD}}\) generations.

\subsection{Evaluation}\label{evaluation}

Because direct transmission events constitute a small fraction of all
possible host pairs, we evaluate inference performance such as IBD, IBS
and phylogenetic patristic distances using precision--recall curves and
the area under the precision--recall curve (AUPRC), which is more
informative than ROC-based metrics under class imbalance. Performance
metrics were computed independently for each simulation replicate. For
each recombination rate, results were summarized across replicates using
the mean and interquartile range. Trends across recombination regimes
were assessed using log-scaled recombination rates.

\subsection{Simulation Workflow}\label{simulation-workflow}

\subsubsection{Stage 1: Forward Simulation with
SLiM}\label{stage-1-forward-simulation-with-slim}

We used SLiM 4.0 to simulate \(G = 200\) generations of forward-time
evolution with the following workflow:

\textbf{Initialization} (generation 1):

\begin{itemize}
\tightlist
\item
  Create single population \(p_1\) with \(N = 10,000\) haploid
  individuals
\item
  Initialize chromosome length \(L = 640,851\) bp (chromosome 1)
\item
  Set recombination rate \(r\)
\item
  Enable tree sequence recording with pedigree tracking
\item
  Define neutral mutation type \(m_1\) and selected mutation type
  \(m_2\)
\end{itemize}

\textbf{Population size dynamics} (generations 1 to
\(200 - g_{\text{change}}\)):

- Implement exponential growth/decline according to \(N(t)\) equation

- For generation \(t\):

\[N_t = N_0 \left(\frac{N_e}{N_0}\right)^{(200-t)/g_{\text{change}}}\]

\textbf{Selection initiation} (generation \(200 - g_s - 1 = 119\)):

- Save simulation state as checkpoint

- Introduce beneficial mutation at position \(\ell_s\) in
\(n_{\text{origins}}\) random genomes

- Track derived allele frequency (DAF) each generation

\textbf{Conditioning on establishment} (generations \(119\) to \(200\)):

- Monitor derived allele frequency \(p_t\) each generation

- If \(p_t = 0\) (loss) and restart count \(< 100\):

- Reload checkpoint state - Re-introduce mutation with new random seed

- Increment restart counter - If \(p_t = 1\) (fixation): - End tracking
(mutation has fixed) - Continue until generation \(200\)

\textbf{High relatedness simulation} (optional, final 50 generations):

- If \texttt{sim\_relatedness\ =\ 1}, activate \texttt{modifyChild()}
callback

- Filter mating pairs to prefer relatedness \(1/32 < r_{ij} < 1/4\)

- This enriches for specific relationship classes (half-siblings,
cousins)

- Used to evaluate IBD methods on known pedigree relationships

\textbf{Output}:

- Tree sequence file (\texttt{.trees}) with complete genealogical
history

- Derived allele frequency trajectory

- True effective population size (\(N_e\)) trajectory

- Restart count (number of failed establishment attempts)

\subsubsection{Stage 2: Sampling and
Simplification}\label{stage-2-sampling-and-simplification}

From the final generation (\(t = 200\)), we sampled individuals to
represent genomic surveillance:

\textbf{Sample selection}: - Identify all individuals alive at \(t = 0\)
(present day): \(\mathcal{I}_{\text{alive}}\) - Randomly sample
\(n_{\text{sample}}/2\) individuals from \(\mathcal{I}_{\text{alive}}\)
- Division by 2 because each haploid individual has 2 nodes in tree
sequence representation - Extract node IDs for sampled individuals - For
multi-population simulations, sample proportionally from each deme.

We varied sampling intensity from Low: 5\% of outbreak
(\(n_{\text{sample}}/N_{\text{outbreak}} = 0.05\)), Medium: 10-20\% of
outbreak and High: 30\% of outbreak (approaching exhaustive sampling).

\textbf{Tree sequence simplification}: We Reduce tree sequence to
sampled lineages: \(\texttt{ts.simplify}(\text{sample\_nodes})\), remove
non-ancestral branches (grafting lineages that don't coalesce to
samples) and retain input roots to preserve pre-simulation history. This
reduces file size and computational burden for downstream analyses

\subsubsection{Stage 3: Recapitation and Mutation
Addition}\label{stage-3-recapitation-and-mutation-addition}

Forward simulation in SLiM captures recent evolutionary dynamics but
does not extend to complete coalescence (infinite time). We used msprime
to ``recapitate'' the tree sequence, adding ancient genealogical
history. This extends all uncoalesced lineages backward in time under a
neutral Wright-Fisher model until complete coalescence, ensuring every
genomic position has a most recent common ancestor (MRCA). After
recapitation, we removed all SLiM-generated mutations (which were used
only for selection tracking) and added neutral mutations uniformly
across the genome. Mutations are placed on tree branches according to a
Poisson process with rate \(\mu\), generating realistic patterns of
standing variation and derived polymorphisms. This allowed us to
generate a recapitated tree sequence, mutated tree sequence with final
variant set and VCF file with pseudo-homozygous genotypes used for the
downstream analyses.

\subsection{Parameter Space
Exploration}\label{parameter-space-exploration}

We simulated three nested experimental designs:

\subsubsection{Scenario 1: Baseline (Fixed
Parameters)}\label{scenario-1-baseline-fixed-parameters}

Single replicate with P. \textbf{\emph{falciparum}} typical parameters:
- \(N_e = 10,000\), \(N_0 = 1,000\) - \(r = 6.67 \times 10^{-7}\)
bp\(^{-1}\) gen\(^{-1}\) (intermediate recombination) -
\(\mu = 10^{-8}\) bp\(^{-1}\) gen\(^{-1}\) - \(B = 5\) (moderate
bottleneck) - \(\sigma = 0.95\) (high selfing) - \(s = 0.3\),
\(g_s = 80\) - \(n_{\text{sample}} = 100\) - Single population (no
migration)

\textbf{Purpose}: Establish baseline accuracy of IBD/IBS/phylogenetic
patristic methods under realistic conditions.

\subsubsection{Scenario 2: Recombination
Sweep}\label{scenario-2-recombination-sweep}

Systematic variation of recombination rate while holding other
parameters constant: - \(r \in \{10^{-9}, 10^{-8}, 10^{-7}, 10^{-6}\}\)
bp\(^{-1}\) gen\(^{-1}\) - All other parameters as in Scenario 1 - 10
replicates per recombination rate - Total: \(4 \times 10 = 40\)
simulations

\textbf{Purpose}: Isolate the effect of recombination on transmission
inference accuracy.

\subsubsection{Scenario 3: Full Parameter
Space}\label{scenario-3-full-parameter-space}

Factorial design crossing:

\begin{itemize}
\tightlist
\item
  \textbf{Recombination rate}:
  \(r \in \{10^{-9}, 10^{-8}, 10^{-7}, 10^{-6}\}\) (4 levels)
\item
  \textbf{Bottleneck size}: \(B \in \{1, 5, 20\}\) (3 levels)
\item
  \textbf{Expected substitutions/transmission}: EST
  \(\in \{0, 0.1, 0.5\}\) (3 levels)
\item
  \textbf{Sampling proportion}: \(f \in \{0.05, 0.10, 0.20\}\) (3
  levels)
\item
  \textbf{Outbreak size}: \(N_{\text{outbreak}} = 2000\) (fixed)
\item
  \textbf{Number of populations}: \(n_{\text{pop}} = 3\) (fixed)
\item
  \textbf{Migration rate}: \(m \in \{0.001, 0.01, 0.05\}\) (3 levels)
\end{itemize}

Core biological scenarios: \(4 \times 3 \times 3 = 36\) combinations
With sampling/migration variation: \(36 \times 3 \times 3 = 324\)
combinations With 20 replicates: \(324 \times 20 = 6,480\) total
simulations.

\textbf{Purpose}: Comprehensively map the parameter space where genomic
methods can accurately infer transmission, identifying critical
thresholds and interaction effects.

\subsection{Output and Data
Availability}\label{output-and-data-availability}

For each simulation, we generated:

\textbf{Primary outputs}: - Tree sequence files (\texttt{.trees}):
Complete ARG with metadata - VCF files (\texttt{.vcf.gz}): Biallelic
SNPs with phased haploid genotypes - Ground truth transmission network:
Known donor-recipient pairs from pedigree - Parameter files
(\texttt{.json}): Complete simulation metadata

\textbf{Diagnostic outputs}: - SLiM stdout logs (restart counts,
population size trajectories) - Recapitation success metrics - Mutation
count summaries

\subsection{Empirical dataset for distributional consistency
validation}\label{empirical-dataset-for-distributional-consistency-validation}

To assess whether the simulated performance bounds are consistent with
empirical P. \textbf{falciparum} genomic data, we applied the three
inference methods to whole-genome sequence data from 160 P.
\textbf{falciparum} isolates collected i 2014 from ten villages in The
Gambia (MRC Unit The Gambia at LSHTM, 2014--2015; Table SX). Samples
were collected from households enrolled in a longitudinal malaria
surveillance programme. Household membership (HHCode, VillageCode) and
visit date were used to define epidemiological proximity. Because
confirmed transmission ground truth is unavailable in this observational
dataset, same-household pairs sampled within the same transmission
season were used as a proxy for transmission-linked pairs, following the
approach of Daniels et al.~(2015) and Taylor et al.~(2017). IBD
inference was restricted to monoclonal infections (COIL = 1; n = 152) to
satisfy hmmIBD input assumptions. All three methods were applied using
the same pipeline as the simulation analyses, and AUPRC was computed
using household proximity as the outcome label.

All simulation codes, parameter files, and analysis scripts are
available at {[}GitHub repository URL{]}. Simulated datasets are
archived at {[}data repository DOI{]}.

\newpage

\section*{References}\label{references}
\addcontentsline{toc}{section}{References}

\phantomsection\label{refs}
\begin{CSLReferences}{1}{0}
\bibitem[\citeproctext]{ref-Albers2020}
Albers, Patrick K., and Gil McVean. 2020. {``Dating Genomic Variants and
Shared Ancestry in Population-Scale Sequencing Data.''} \emph{PLoS
Biology} 18 (1): e3000586.

\bibitem[\citeproctext]{ref-AmambuaNgwa2019}
Amambua-Ngwa, Alfred, Lucas Amenga-Etego, Edwin Kamau, Roberto Amato,
Anita Ghansah, Lemu Golassa, Milijaona Randrianarivelojosia, et al.
2019. {``Major Subpopulations of Plasmodium Falciparum in Sub-Saharan
Africa.''} \emph{Science} 365 (6455): 813--16.
\url{https://doi.org/10.1126/science.aav5427}.

\bibitem[\citeproctext]{ref-Apinjoh2019}
Apinjoh, Tobias O., Amed Ouattara, Vincent P. K. Titanji, Abdoulaye
Djimde, and Alfred Amambua-Ngwa. 2019. {``Genetic Diversity and Drug
Resistance Surveillance of Plasmodium Falciparum for Malaria
Elimination: Is There an Ideal Tool for Resource-Limited Sub-Saharan
Africa?''} \emph{Malaria Journal} 18 (1).
\url{https://doi.org/10.1186/s12936-019-2844-5}.

\bibitem[\citeproctext]{ref-Auburn2018}
Auburn, Sarah, Ernest D. Benavente, Olivo Miotto, Richard D. Pearson,
Roberto Amato, Matthew J. Grigg, Bridget E. Barber, et al. 2018.
{``Genomic Analysis of a Pre-Elimination Malaysian Plasmodium Vivax
Population Reveals Selective Pressures and Changing Transmission
Dynamics.''} \emph{Nature Communications} 9 (1).
\url{https://doi.org/10.1038/s41467-018-04965-4}.

\bibitem[\citeproctext]{ref-baumdicker2022}
Baumdicker, Franz, Gertjan Bisschop, Daniel Goldstein, Graham Gower,
Aaron P Ragsdale, Georgia Tsambos, Sha Zhu, et al. 2022. {``Efficient
Ancestry and Mutation Simulation with Msprime 1.0.''} \emph{Genetics}
220 (3): iyab229.

\bibitem[\citeproctext]{ref-Campbell2018}
Campbell, Fraser, Anne Cori, Neil Ferguson, and Thibaut Jombart. 2018.
{``Bayesian Inference of Transmission Chains Using Timing of Symptoms,
Pathogen Genomes and Contact Data.''} \emph{PLoS Computational Biology}
14 (4): e1006205.

\bibitem[\citeproctext]{ref-chenet2012}
Chenet, Stella M, Kristan A Schneider, Leopoldo Villegas, and Ananias A
Escalante. 2012. {``Local Population Structure of Plasmodium: Impact on
Malaria Control and Elimination.''} \emph{Malaria Journal} 11 (1): 412.

\bibitem[\citeproctext]{ref-Conway1999}
Conway, David J., Cally Roper, Adebola M. J. Oduola, David E. Arnot,
Peter G. Kremsner, Martin P. Grobusch, Christopher F. Curtis, and Brian
M. Greenwood. 1999. {``High Recombination Rate in Natural Populations of
Plasmodium Falciparum.''} \emph{Proceedings of the National Academy of
Sciences} 96 (8): 4506--11.

\bibitem[\citeproctext]{ref-Daniels2015}
Daniels, Rachel F., Stephen F. Schaffner, Edward A. Wenger, Joshua L.
Proctor, Hsiao-Han Chang, Willie Wong, Nathalie Baro, et al. 2015.
{``Modeling Malaria Genomics Reveals Transmission Decline and Rebound in
Senegal.''} \emph{Proceedings of the National Academy of Sciences} 112
(22): 7067--72.

\bibitem[\citeproctext]{ref-DeMaio2015}
De Maio, Nicola, Chieh-Hsi Wu, Kathleen M O'Reilly, and Daniel Wilson.
2015. {``New Routes to Phylogeography: A Bayesian Structured Coalescent
Approximation.''} Edited by Jonathan K. Pritchard. \emph{PLOS Genetics}
11 (8): e1005421. \url{https://doi.org/10.1371/journal.pgen.1005421}.

\bibitem[\citeproctext]{ref-Didelot2017}
Didelot, Xavier, Christophe Fraser, Jennifer Gardy, and Caroline Colijn.
2017. {``Genomic Infectious Disease Epidemiology in Partially Sampled
and Ongoing Outbreaks.''} \emph{Molecular Biology and Evolution} 34 (4):
997--1007.

\bibitem[\citeproctext]{ref-Didelot2014}
Didelot, Xavier, Jennifer Gardy, and Caroline Colijn. 2014. {``Bayesian
Inference of Infectious Disease Transmission from Whole-Genome Sequence
Data.''} \emph{Molecular Biology and Evolution} 31 (7): 1869--79.
\url{https://doi.org/10.1093/molbev/msu121}.

\bibitem[\citeproctext]{ref-Didelot2021}
Didelot, Xavier, Mark Kendall, Yinyin Xu, Sarah White, and Anne Cori.
2021. {``Genomic Epidemiology Analysis of Infectious Disease Outbreaks
Using Dated Phylogenies.''} \emph{Trends in Microbiology} 29 (6):
476--86.

\bibitem[\citeproctext]{ref-Didelot2020}
Didelot, Xavier, Igor Siveroni, and Erik M Volz. 2020. {``Additive
Uncorrelated Relaxed Clock Models for the Dating of Genomic Epidemiology
Phylogenies.''} Edited by Keith Crandall. \emph{Molecular Biology and
Evolution} 38 (1): 307--17.
\url{https://doi.org/10.1093/molbev/msaa193}.

\bibitem[\citeproctext]{ref-DiezBenavente2020}
Diez Benavente, Ernest, Monica Campos, Jody Phelan, Debbie Nolder,
Jamille G. Dombrowski, Claudio R. F. Marinho, Kanlaya Sriprawat, et al.
2020. {``A Molecular Barcode to Inform the Geographical Origin and
Transmission Dynamics of Plasmodium Vivax Malaria.''} Edited by Giorgio
Sirugo. \emph{PLOS Genetics} 16 (2): e1008576.
\url{https://doi.org/10.1371/journal.pgen.1008576}.

\bibitem[\citeproctext]{ref-Felger2012}
Felger, Ingrid, Martin Maire, Michael T. Bretscher, Nicole Falk, André
Tiaden, Wilson Sama, Hans-Peter Beck, Seth Owusu-Agyei, and Thomas A.
Smith. 2012. {``The Dynamics of Natural Plasmodium Falciparum
Infections.''} Edited by Roly D. Gosling. \emph{PLoS ONE} 7 (9): e45542.
\url{https://doi.org/10.1371/journal.pone.0045542}.

\bibitem[\citeproctext]{ref-ferreira2007}
Ferreira, Marcelo U, Nadira D Karunaweera, Mônica da Silva-Nunes, Natal
S Da Silva, Dyann F Wirth, and Daniel L Hartl. 2007. {``Population
Structure and Transmission Dynamics of Plasmodium Vivax in Rural
Amazonia.''} \emph{The Journal of Infectious Diseases} 195 (8):
1218--26.

\bibitem[\citeproctext]{ref-Griffin2010}
Griffin, Jamie T., T. Déirdre Hollingsworth, Lucy C. Okell, Thomas S.
Churcher, Michael White, W. Hinsley, Teun Bousema, Chris Drakeley, Neil
M. Ferguson, and Azra C. Ghani. 2010. {``Reducing Plasmodium Falciparum
Malaria Transmission in Africa: A Model-Based Evaluation of Intervention
Strategies.''} \emph{PLoS Medicine} 7 (8): e1000324.

\bibitem[\citeproctext]{ref-haller2019}
Haller, Benjamin C, and Philipp W Messer. 2019. {``SLiM 3: Forward
Genetic Simulations Beyond the Wright--Fisher Model.''} \emph{Molecular
Biology and Evolution} 36 (3): 632--37.

\bibitem[\citeproctext]{ref-Hay2002}
Hay, Simon I., Milka Simba, Millie Busolo, Abdisalan M. Noor, Helen L.
Guyatt, Sam A. Ochola, and Robert W. Snow. 2002. {``Defining and
Detecting Malaria Epidemics in the Highlands of Western Kenya.''}
\emph{Emerging Infectious Diseases} 8 (6): 555--62.
\url{https://doi.org/10.3201/eid0806.010310}.

\bibitem[\citeproctext]{ref-Henden2018}
Henden, Lyndal, Stuart Lee, Ivo Mueller, Alyssa Barry, and Melanie
Bahlo. 2018. {``Identity-by-Descent Analyses for Measuring Population
Dynamics and Selection in Recombining Pathogens.''} Edited by Xavier
Didelot. \emph{PLOS Genetics} 14 (5): e1007279.
\url{https://doi.org/10.1371/journal.pgen.1007279}.

\bibitem[\citeproctext]{ref-Jiang2011}
Jiang, Hongying, Na Li, Vivek Gopalan, Martine M Zilversmit, Sudhir
Varma, Vijayaraj Nagarajan, Jian Li, et al. 2011. {``High Recombination
Rates and Hotspots in a Plasmodium Falciparum Genetic Cross.''}
\emph{Genome Biology} 12 (4).
\url{https://doi.org/10.1186/gb-2011-12-4-r33}.

\bibitem[\citeproctext]{ref-Kelleher2018}
Kelleher, Jerome, Kevin R. Thornton, Jaime Ashander, and Peter L. Ralph.
2018. {``Efficient Pedigree Recording for Fast Population Genetics
Simulation.''} \emph{PLoS Computational Biology} 14 (11): e1006581.

\bibitem[\citeproctext]{ref-MalariaGEN2023}
MalariaGEN, Ambroise Ahouidi, Mozam Ali, et al. 2021. {``An Open Dataset
of {Plasmodium falciparum} Genome Variation in 7,000 Samples from 28
Countries.''} \emph{Wellcome Open Research} 6: 42.
\url{https://doi.org/10.12688/wellcomeopenres.16168.2}.

\bibitem[\citeproctext]{ref-Manske2012}
Manske, Magnus, Olivo Miotto, Susana Campino, Sarah Auburn, Jacob
Almagro-Garcia, Gareth Maslen, Jack O'Brien, et al. 2012. {``Analysis of
Plasmodium Falciparum Diversity in Natural Infections by Deep
Sequencing.''} \emph{Nature} 487 (7407): 375--79.
\url{https://doi.org/10.1038/nature11174}.

\bibitem[\citeproctext]{ref-Miles2016}
Miles, Alistair, Zamin Iqbal, Paul Vauterin, Richard D. Pearson, Susana
Campino, Michel Theron, Katherine Gould, et al. 2016. {``Indels,
Structural Variation, and Recombination Drive Genomic Diversity in
Plasmodium Falciparum.''} \emph{Genome Research} 26 (9): 1288--99.

\bibitem[\citeproctext]{ref-MohdAbdRazak2016}
Mohd Abd Razak, Mohd Ridzuan, Umi Rubiah Sastu, Nor Azrina Norahmad,
Abass Abdul-Karim, Amirrudin Muhammad, Prem Kumar Muniandy, Jenarun
Jelip, et al. 2016. {``Genetic Diversity of Plasmodium Falciparum
Populations in Malaria Declining Areas of Sabah, East Malaysia.''}
Edited by Georges Snounou. \emph{PLOS ONE} 11 (3): e0152415.
\url{https://doi.org/10.1371/journal.pone.0152415}.

\bibitem[\citeproctext]{ref-Neafsey2015a}
Neafsey, Daniel E., Michal Juraska, Trevor Bedford, David Benkeser,
Clarissa Valim, Allison Griggs, Marc Lievens, et al. 2015. {``Genetic
Diversity and Protective Efficacy of the RTS, s/AS01 Malaria Vaccine.''}
\emph{New England Journal of Medicine} 373 (21): 2025--37.
\url{https://doi.org/10.1056/nejmoa1505819}.

\bibitem[\citeproctext]{ref-Neafsey2015b}
Neafsey, Daniel E., Robert M. Waterhouse, Mohammad R. Abai, Sergey S.
Aganezov, Max A. Alekseyev, James E. Allen, James Amon, et al. 2015.
{``Highly Evolvable Malaria Vectors: The Genomes of 16 Anopheles
Mosquitoes.''} \emph{Science} 347 (6217).
\url{https://doi.org/10.1126/science.1258522}.

\bibitem[\citeproctext]{ref-Nkhoma2012}
Nkhoma, Standwell C., Shalini Nair, Ian H. Cheeseman, Cherise
Rohr-Allegrini, Sittaporn Singlam, François Nosten, and Tim J. C.
Anderson. 2012. {``Close Kinship Within Multiple-Genotype Malaria
Parasite Infections.''} \emph{Proceedings of the Royal Society B:
Biological Sciences} 279 (1738): 2589--98.
\url{https://doi.org/10.1098/rspb.2012.0113}.

\bibitem[\citeproctext]{ref-Nkhoma2020}
Nkhoma, Standwell C., Simon G. Trevino, Karla M. Gorena, Shalini Nair,
Stanley Khoswe, Catherine Jett, Roy Garcia, et al. 2020.
{``Co-Transmission of Related Malaria Parasite Lineages Shapes
Within-Host Parasite Diversity.''} \emph{Cell Host \&Amp; Microbe} 27
(1): 93--103.e4. \url{https://doi.org/10.1016/j.chom.2019.12.001}.

\bibitem[\citeproctext]{ref-schaffner2018}
Schaffner, Stephen F, Aimee R Taylor, Wesley Wong, Dyann F Wirth, and
Daniel E Neafsey. 2018. {``hmmIBD: Software to Infer Pairwise Identity
by Descent Between Haploid Genotypes.''} \emph{Malaria Journal} 17 (1):
196.

\bibitem[\citeproctext]{ref-smith2010}
Smith, David L, Chris J Drakeley, Christinah Chiyaka, and Simon I Hay.
2010. {``A Quantitative Analysis of Transmission Efficiency Versus
Intensity for Malaria.''} \emph{Nature Communications} 1 (1): 108.

\bibitem[\citeproctext]{ref-Speed2015}
Speed, Doug, and David J. Balding. 2015. {``Relatedness in the
Post-Genomic Era: Is It Still Useful?''} \emph{Nature Reviews Genetics}
16 (1): 33--44.

\bibitem[\citeproctext]{ref-taylor2017}
Taylor, Aimee R, Stephen F Schaffner, Gustavo C Cerqueira, Standwell C
Nkhoma, Timothy JC Anderson, Kanlaya Sriprawat, Aung Pyae Phyo, François
Nosten, Daniel E Neafsey, and Caroline O Buckee. 2017. {``Quantifying
Connectivity Between Local Plasmodium Falciparum Malaria Parasite
Populations Using Identity by Descent.''} \emph{PLoS Genetics} 13 (10):
e1007065.

\bibitem[\citeproctext]{ref-Taylor2017}
Taylor, Alyssa R., Pradipsinh E. Jacob, Daniel E. Neafsey, and Caroline
O. Buckee. 2017. {``Quantifying Connectivity Between Local Plasmodium
Falciparum Malaria Parasite Populations Using Identity by Descent.''}
\emph{PLoS Genetics} 13 (10): e1007065.

\bibitem[\citeproctext]{ref-Wesolowski2018}
Wesolowski, Amy, Alyssa R. Taylor, Hsiao-Han Chang, Robert Verity,
Sachel Mok Tessema, Jeffrey A. Bailey, and Caroline O. Buckee. 2018.
{``Mapping Malaria by Combining Parasite Genomic and Epidemiological
Data.''} \emph{Science} 362 (6419): eaar6872.

\bibitem[\citeproctext]{ref-Wong2025}
Wong, Thomas, Nhan Ly-Trong, Huaiyan Ren, Hector Baños, Andrew Roger,
Edward Susko, Chris Bielow, et al. 2025. {``IQ-TREE 3: Phylogenomic
Inference Software Using Complex Evolutionary Models,''} April.
\url{https://doi.org/10.32942/x2p62n}.

\bibitem[\citeproctext]{ref-wong2017}
Wong, Wesley, Allison D Griggs, Rachel F Daniels, Stephen F Schaffner,
Daouda Ndiaye, Amy K Bei, Awa B Deme, et al. 2017. {``Genetic
Relatedness Analysis Reveals the Cotransmission of Genetically Related
Plasmodium Falciparum Parasites in Thi{è}s, Senegal.''} \emph{Genome
Medicine} 9 (1): 5.

\bibitem[\citeproctext]{ref-WHO2023}
World Health Organization. 2023. \emph{World Malaria Report 2023}.
Geneva, Switzerland: World Health Organization.

\bibitem[\citeproctext]{ref-Zhu2022}
Zhu, Sijia J., Jacob Almagro-Garcia, Gil McVean, et al. 2022. {``The
Origins and Relatedness Structure of Plasmodium Falciparum
Populations.''} \emph{Nature Communications} 13: 5459.

\end{CSLReferences}

Haller BC, Messer PW. 2023. SLiM 4: Multispecies eco-evolutionary
modeling. \emph{American Naturalist} 201(5):E127-E139.

Kelleher J, Etheridge AM, McVean G. 2016. Efficient coalescent
simulation and genealogical analysis for large sample sizes. \emph{PLoS
Computational Biology} 12(5):e1004842.

Hamilton WL, et al.~2017. Evolution and expansion of multidrug-resistant
malaria in southeast Asia: a genomic epidemiology study. \emph{Lancet
Infectious Diseases} 17(5):491-497.

Daniels R, et al.~2015. Modeling malaria genomics reveals transmission
decline and rebound in Senegal. \emph{PNAS} 112(22):7067-7072.

Schneider KA, et al.~2018. Virulence, drug sensitivity and transmission
success in the rodent malaria, Plasmodium \emph{chabaudi}. Proceedings
of the Royal Society B 285:20181391.

\newpage

\subsubsection{Supplementary
Information}\label{supplementary-information}

\subsubsection{Acknowledgements}\label{acknowledgements}

\subsubsection{Author contributions}\label{author-contributions}

\subsubsection{Funding}\label{funding}

\subsubsection{Data availability}\label{data-availability}

\subsubsection{Ethics approval and consent to
participate}\label{ethics-approval-and-consent-to-participate}

\subsubsection{Competing interests}\label{competing-interests}

The authors declare no competing interests

\newpage

\subsection{Supplementary Figures}\label{supplementary-figures}

\pandocbounded{\includegraphics[keepaspectratio]{../results/supplementary/FigS1_baseline_ibd.pdf}}

\textbf{Figure S1}. True pairwise IBD fraction distribution (panel A)
and hmmIBD-inferred IBD fraction distribution (panel B) under the
baseline scenario. Both panels show the full count distribution across
all simulated pairs. The bimodal structure in panel A reflects the
high-selfing, low-recombination baseline parameterization.

\pandocbounded{\includegraphics[keepaspectratio]{../results/supplementary/figureSx_method_comparison.pdf}}

\textbf{Figure Sx}. Sensitivity (true positive rate; panel A) and false
discovery rate (FDR = FP / (TP + FP); panel B) for IBD, IBS, and Phylo
methods across the baseline, recombination sweep, and migration
scenarios. Each dot and whisker represents the median and interquartile
range across replicates. Results shown at a fixed classification
threshold.

\pandocbounded{\includegraphics[keepaspectratio]{../results/supplementary/FigS3_M1_sensitivity_90spec_migration.pdf}}

\textbf{Figure S3}. Sensitivity at 90\% specificity (± 95\% CI) for IBD
(orange), IBS (blue), and Phylo (green) across four recombination rates
(panels) and three migration rates (x-axis) for G ≤ 25 generations.

\pandocbounded{\includegraphics[keepaspectratio]{../results/supplementary/FigS6_P1_migration_auprc_gradient.pdf}}

\textbf{Figure S6}. Mean AUPRC (± 95\% CI) for IBD, IBS, and Phylo
across recombination rates (x-axis, log scale) for three migration rates
(panels: 0.001, 0.01, 0.05), at G ≤ 25 generations. A horizontal dashed
reference line marks AUPRC = 0.6.

\pandocbounded{\includegraphics[keepaspectratio]{../results/supplementary/FigSxxx_migration_sensitivity.pdf}}

\textbf{Figure Sxxx (migration sensitivity)}. Sensitivity at 90\%
specificity (± 95\% CI) as a function of migration rate (x-axis: 0.001,
0.01, 0.05) for each recombination rate stratum (panels A--D: r =
\(10^{-9}\), \(10^{-8}\), \(10^{-7}\), \(10^{-6}\)). Colours indicate
method (IBD, IBS, Phylo).

\pandocbounded{\includegraphics[keepaspectratio]{../results/supplementary/FigS7_BONUS_G_recrate_interaction.pdf}}

\textbf{Figure S7}. Full G-threshold × recombination rate AUPRC heatmaps
for IBD, IBS, and Phylo methods (panels), mirroring the structure of
Fig. 8A--C. Shown separately as a supplementary reference with extended
numerical annotations.

\pandocbounded{\includegraphics[keepaspectratio]{../results/supplementary/FigS2_P3_bottleneck_auprc.pdf}}

\textbf{Figure S2}. Effect of effective population size at the
bottleneck (Ne bottleneck = 1,000; 5,000; 20,000) on mean AUPRC (± 95\%
CI) for IBD, IBS, and Phylo, shown separately for low (r = \(10^{-9}\))
and medium (r = \(10^{-8}\)) recombination rates. Results are for G ≤ 25
generations.

\pandocbounded{\includegraphics[keepaspectratio]{../results/supplementary/FigS4_M2_identifiable_fraction_migration.pdf}}

\textbf{Figure S4}. Fraction of simulation replicates classified as
identifiable (i.e., achieving high-performance AUPRC above a fixed
threshold) for IBD, IBS, and Phylo methods, across three migration rates
(rows) and four recombination rates (columns). For G ≤ 25 generations.
Cells show the fraction (0--1); colour scale from dark red (fraction =
0) to dark green (fraction = 1).

\pandocbounded{\includegraphics[keepaspectratio]{../results/supplementary/FigS5_P4_ibd_method_gap_heatmap.pdf}}

\textbf{Figure S5}. IBD advantage over the next-best method (AUPRC gap =
AUPRC\_IBD − max(AUPRC\_IBS, AUPRC\_Phylo)) for G ≤ 25 generations.
Heatmap with migration rate on the y-axis (0.001, 0.01, 0.05) and
recombination rate on the x-axis (\(10^{-9}\), \(10^{-8}\), \(10^{-7}\),
\(10^{-6}\)). Numeric values shown in each cell; colour scale from green
(small gap) to dark red (large gap).

\newpage

\subsection{Potential Tables for the
Manuscript}\label{potential-tables-for-the-manuscript}

\textbf{Table 1. Simulation parameters and scenario definitions} A
structured summary of all three scenarios (Baseline, Recombination
Sweep, Multi-population), listing the values of μ, r, Ne, selfing rate,
drug selection coefficient, number of subpopulations, migration rates,
and the factorial combinations explored. This table contextualises all
downstream results and is essential for reproducibility.

\textbf{Table 2. AUPRC summary by method, scenario, and G threshold
(main conditions)} Mean AUPRC (± 95\% CI) for IBD, IBS, and Phylo at G =
1, 5, 10, and 25, for each scenario and key recombination rate (e.g., r
= \(10^{-9}\) and r = \(10^{-6}\)). This table complements Fig. 8 and
provides exact values for readers wishing to compare performance at
specific conditions.

\textbf{Table 3. Sensitivity at 90\% specificity across recombination
and migration (key conditions)} Mean sensitivity at 90\% specificity (±
95\% CI) for IBD, IBS, and Phylo at each recombination × migration
combination for G ≤ 25 generations. This provides a fixed-specificity
operating point complement to the AUPRC-based summaries, facilitating
comparison with threshold-based field application scenarios. Complements
Fig. S3 and FigSxxx\_migration\_sensitivity.

\textbf{Table 4. Observed versus simulated AUPRC and matched simulation
conditions} For each method (IBD, IBS, Phylo), the observed AUPRC in the
real P. \emph{falciparum} dataset, the range of simulated AUPRC values
across all conditions, and the simulated condition (recombination rate,
migration, sample size) that most closely matches the observed
performance. Complements Fig. 9B and anchors the biological
interpretation of the simulation results.

\end{document}
